\documentclass{article}

% The root Makefile supplies the path to the unmodified ICML style files.
% Shared, topic-agnostic preamble for all three papers.
% Keep formatting changes here conservative: ICML forbids modifications that
% create a space advantage.
\usepackage{microtype}
\usepackage{graphicx}
\usepackage{subcaption}
\usepackage{booktabs}
\usepackage{hyperref}

% Improve interoperability between hyperref and algorithmic.
\newcommand{\theHalgorithm}{\arabic{algorithm}}

% Review mode is intentionally the default. Change this to
% \usepackage[accepted]{icml2026} only after acceptance.
\usepackage{icml2026}

\usepackage{amsmath}
\usepackage{amssymb}
\usepackage{mathtools}
\usepackage{amsthm}
\usepackage[capitalize,noabbrev]{cleveref}

\theoremstyle{plain}
\newtheorem{theorem}{Theorem}[section]
\newtheorem{proposition}[theorem]{Proposition}
\newtheorem{lemma}[theorem]{Lemma}
\newtheorem{corollary}[theorem]{Corollary}
\theoremstyle{definition}
\newtheorem{definition}[theorem]{Definition}
\newtheorem{assumption}[theorem]{Assumption}
\theoremstyle{remark}
\newtheorem{remark}[theorem]{Remark}

\graphicspath{{figures/}}

% Visible by design: no template should be mistaken for a completed paper.
\newcommand{\placeholder}[1]{%
  \par\noindent\textit{Template note: #1}\par%
}


\icmltitlerunning{Decision-Calibrated Paired Evaluation}

\begin{document}

\twocolumn[
  \icmltitle{Decision-Calibrated Paired Evaluation with
  Sequentially Acquired Human Judgments}

  % Author details remain hidden in review mode. Replace these placeholders
  % only when preparing the camera-ready version.
  \begin{icmlauthorlist}
    \icmlauthor{Anonymous Author(s)}{anon}
  \end{icmlauthorlist}
  \icmlaffiliation{anon}{Anonymous Institution}
  \icmlcorrespondingauthor{Anonymous Author}{anonymous@example.com}

  \icmlkeywords{language-model evaluation, human preference, adaptive
  sampling, uncertainty quantification, decision making}

  \vskip 0.3in
]

\printAffiliationsAndNotice{}

\begin{abstract}
Pairwise language-model evaluations increasingly use inexpensive model judges,
but deployment decisions concern a specified human population and must account
for ties, prompt variation, correlated judge errors, and adaptively chosen
human labels. We formulate decision-calibrated paired evaluation around the
finite-frame estimand ``human win probability minus loss probability'' under
fixed item, prompt, response-generation, rater, and display-order
distributions. A tie-aware hierarchical jury model drives one-step
value-of-information acquisition, while an explicit random-exploration mixture
and logged draw probabilities preserve design-based identification. A
prequential augmented Hansen--Hurwitz estimator is exactly design-unbiased
and supports fixed-budget martingale inference under stated positivity and
boundedness conditions, even when the jury model is misspecified. This
methods-first draft specifies simulation, human-audit, and temporal-shift tests;
efficiency and decision-quality claims remain pending empirical execution.
\end{abstract}

\section{Introduction}
\label{sec:introduction}

Model comparison is often reduced to a leaderboard score, although the actual
decision may be whether system $A$ should replace system $B$ for a particular
population of requests. Human pairwise evaluation directly targets preference
and underlies Chatbot Arena \cite{chiang2024arena}, but comprehensive human
labeling is expensive. Model-based judges scale cheaply
\cite{zheng2023judging} yet exhibit position, style, verbosity, and
self-preference effects \cite{panickssery2024evaluators,ye2025justice}. A jury
of correlated judges cannot be treated as an equally large set of independent
humans.

Recent work estimates judge-specific bias and uncertainty
\cite{dubois2026biases}, shared confounding \cite{zhao2026care}, jury
discrimination and internal consistency \cite{qian2026jury}, and uncertainty
with iterative comparison selection in comparative evaluation
\cite{fathullah2025uncertainty}. Human-calibrated estimators correct noisy
model judgments \cite{lee2026report,chen2026efficient}. Efficient inference for
a static, independently selected calibration set is now well understood enough
that ``apply AIPW to judge labels'' is not a new contribution.

The unresolved issue studied here is sequential acquisition. Once human cells
are chosen because previous human outcomes made them decision-relevant, the
calibration set is not missing completely at random. A final inclusion
probability cannot generally be reconstructed by multiplying probabilities
along the realized outcome path. Ordinary post-hoc cross-fitting can also leak
future selection information. Meanwhile, pure uncertainty or value-of-
information acquisition can assign zero probability to parts of the target
population, breaking design-based inference.

We ask: \emph{can adaptive human labeling reduce the cost of deciding whether
$A$ is materially better than $B$ while retaining valid inference for an
explicit human-preference estimand?} We make four scoped contributions:
\begin{itemize}
  \item We define a tie-aware finite evaluation frame that fixes item, prompt,
  response coupling, rater, and display-order distributions before labeling.
  \item We use a crossed jury model only for prediction and decision-targeted
  acquisition, not as the definition of quality.
  \item We combine value-of-information acquisition with randomized exploration
  and a draw-level sequential augmented Hansen--Hurwitz estimator. We state
  its exact design-unbiasedness and fixed-budget asymptotic conditions.
  \item We pre-specify simulations, a sealed target-random human audit, shift tests,
  and ablations that separate interval validity from label efficiency.
\end{itemize}

\paragraph{Evidence status.}
The formal design claims are established below. Whether adaptive acquisition
saves labels, whether the working model predicts humans, and whether the chosen
human panel represents deployment stakeholders remain empirical questions; no
such result is reported in this draft.

% Camera-ready only: if the work evaluates technology developed by an author's
% employer or has another substantive financial conflict, place the required
% "Conflict of Interest Disclosure" paragraph at the end of this section.

\section{Related Work}
\label{sec:related}

\paragraph{Measurement variation and judge bias.}
Multi-prompt evaluation models performance distributions across prompt forms
\cite{polo2024multiprompt}. Human studies can be assessed by the probability
that a repeat study preserves a ranking \cite{riley2024replicable}. Model
judges introduce additional structure: self-recognition can produce
self-preference \cite{panickssery2024evaluators}, controlled perturbations
reveal multiple biases \cite{ye2025justice}, and modern hierarchical methods
estimate grader effects, confounding, or discrimination
\cite{dubois2026biases,zhao2026care,qian2026jury}. Our jury model borrows this
measurement perspective. Its novelty is not a new bias taxonomy or latent
quality model.

\paragraph{Human-calibrated inference.}
Binary sensitivity/specificity correction with propagated calibration
uncertainty and pilot-based allocation across true-label classes is studied by
\citet{lee2026report}. \citet{chen2026efficient} unify measurement-error and
prediction-powered estimators for mean human outcomes, including pairwise win
rates, under an independent MCAR calibration/test split. Neither selects
evaluation cells using prior human outcomes. Classical unequal-probability
with-replacement estimation \cite{hansen1943sampling} and inference for
adaptively collected data \cite{hadad2021confidence} supply the statistical
foundation. Prediction-Powered Active Testing (PPAT) independently combines
active-testing weights with a prediction control variate
\cite{caron2026ppat}. Our distinction is decision-loss acquisition for
three-way human comparisons over a crossed evaluation frame, not adaptive
residualization or its central limit theorem by itself.

\paragraph{Adaptive and dynamic evaluation.}
Active acquisition can exploit historical dependence among benchmark items
\cite{li2025active}, while ProEval uses surrogate models for performance
estimation and failure discovery \cite{huang2026proeval}. Such compression can
fail under source--target model shift \cite{zhang2025fewer}. Dynamic benchmarks
reduce static contamination through generated or refreshed tasks
\cite{zhu2024dyval,white2025livebench}; contamination tests can detect some
forms of memorized ordering \cite{oren2024contamination}. Freshness, however,
changes the evaluation distribution and complicates longitudinal claims.
Finally, non-transitive judge preferences can make a scalar rank
baseline-dependent \cite{xu2025nontransitivity}. We therefore target one
prespecified pairwise replacement decision, not a universal leaderboard.

\section{Evaluation Goal and Desiderata}
\label{sec:goal}

At evaluation time $t_0$, freeze models $A,B$, their decoding policies, the
rubric, and a finite frame of $N$ root items. Item $i$ has target weight
$w_i\ge0$, $\sum_iw_i=1$. Cell $c\in\mathcal C_i$ identifies a prompt variant
and paired response realizations with prespecified conditional weight
$a_{c\mid i}$. Independent response draws, shared random seeds, and
deterministic decoding define different estimands because preference is
nonlinear; the coupling is therefore fixed explicitly.

Before acquisition, enroll a finite rater panel
$\mathcal H=\{1,\ldots,H\}$ with nonnegative target weights $\rho_h$ summing to
one. Inference therefore targets this panel; extrapolation to other people is a
separate superpopulation claim. For display order $o\in\{-1,+1\}$, define the
outcome of one standardized rating draw
\begin{equation}
Y_{icho}=
\begin{cases}
+1,&A\text{ preferred},\\
0,&\text{tie},\\
-1,&B\text{ preferred}.
\end{cases}
\end{equation}
and $\mu_{icho}=\mathbb E[Y_{icho}\mid i,c,h,o]$, where the expectation covers
within-rater response variation under the fixed protocol.
The finite-frame target is
\begin{equation}
\Delta_F=\sum_{i=1}^Nw_i\sum_{c\in\mathcal C_i}a_{c\mid i}
\sum_{h\in\mathcal H}\rho_h\frac12\sum_{o\in\{-1,+1\}}\mu_{icho}.
\label{eq:estimand}
\end{equation}
Equivalently, define a target unit $x=(i,c,h,o)\in\mathcal X$ with mass
\begin{equation}
\pi(x)=w_i a_{c\mid i}\rho_h/2,\qquad \sum_{x\in\mathcal X}\pi(x)=1.
\label{eq:target-mass}
\end{equation}
Then $\Delta_F=P_\pi(A\text{ wins})-P_\pi(B\text{ wins})$ and the half-tie win
rate is $(1+\Delta_F)/2$. We separately target
\begin{equation}
\tau_F=\sum_{x\in\mathcal X}\pi(x)P(Y_x=0\mid x)
\end{equation}
because the same $\Delta_F$ can describe decisive disagreement or widespread
ties.

A practical margin $\delta>0$ and loss matrix are pre-specified. The action set
is ``$A$ materially better'' if $\Delta_F>\delta$, ``practically equivalent''
if $|\Delta_F|\le\delta$, ``$B$ materially better'' if
$\Delta_F<-\delta$, and ``not established'' when data do not support a
decisive action. This is an operational human-preference construct, not latent
universal quality.

\section{Evaluation Method}
\label{sec:method}

The protocol first collects inexpensive labels from a family-diverse model
jury on all cells, then a uniformly randomized human pilot. Before each
subsequent batch, a working model is trained only on prior batches and used to
compute decision value. Human acquisition mixes this adaptive policy with the
target item distribution. A sealed, uniformly sampled human audit is never
used for fitting or acquisition.

\subsection{Data and Task Construction}

The frame is a probability sample from a timestamped deployment log or a
weighted benchmark with a clearly narrower finite-frame interpretation.
Sampling excludes personal data, unsupported licenses, and items requiring
professional expertise absent from the rater pool. Root-item provenance,
inclusion weight, domain, language, safety category, near-duplicate cluster,
prompt generator, response seeds, model revisions, tokenizer, and decoding
parameters are versioned before labels are revealed.

Each prompt variant is validated as meaning-preserving by a separate process;
the target weights over variants are fixed, not chosen after seeing which model
benefits. Model identifiers are hidden (although style may reveal identity),
AB/BA display is randomized, and humans
receive a real tie option with a fixed definition. ``Equally good'' and ``both
unacceptable'' are recorded separately even if the primary analysis maps both
to zero. Independent ratings define \cref{eq:estimand}; majority vote and
adjudication are secondary, different estimands.

Rater and order assignments follow the target masses in
\cref{eq:target-mass}, or their logged adaptive counterparts. Availability is
resolved before content is shown. Every assignment attempt is logged; a skip
after viewing content is a separate outcome, never silently reassigned. The
primary preference estimand then requires a stated response assumption or
inverse-response correction, and sensitivity to outcome-dependent skipping is
mandatory.

\subsection{Scoring and Aggregation}

At round $t$, sample one complete target unit with replacement,
\begin{equation}
X_t=(I_t,C_t,H_t,O_t)\sim
p_t(\cdot\mid\mathcal F_{t-1}),
\end{equation}
where $\mathcal F_{t-1}$ contains every prior assignment, sampling
probability, jury feature, outcome, and model random seed. This joint law makes
rater and order assignment explicit. If operational constraints require a
nested assignment, the product of logged item, cell, rater, and order
probabilities must equal $p_t(X_t)$.

Let $m_{t-1}(x)\in[-1,1]$ be a bounded prediction of $Y_x$ trained only on
prior rounds and frozen before draw $t$, and define
\begin{equation}
\bar m_{t-1}=\sum_{x\in\mathcal X}\pi(x)m_{t-1}(x).
\end{equation}
The draw-level sequential augmented Hansen--Hurwitz contribution is
\begin{equation}
\psi_t=\bar m_{t-1}+
\frac{\pi(X_t)}{p_t(X_t)}
\{Y_{X_t}-m_{t-1}(X_t)\},
\label{eq:psi}
\end{equation}
with final estimator
\begin{equation}
\widehat\Delta_T=\frac1T\sum_{t=1}^T\psi_t.
\label{eq:estimator}
\end{equation}
Repeated draws are retained. $p_t$ is the conditional draw probability, not
the probability that a unit is ever included; substituting a realized-path
inclusion probability is generally invalid under outcome adaptation. Replacing
$Y_x$ in \cref{eq:psi} by $\mathbf{1}\{Y_x=0\}$ gives the tie-rate estimator,
with its covariance with $\widehat\Delta_T$ reported.

\subsection{Human or Model-Based Judging}

The jury model is a prediction and diagnostic instrument, not the source of
identification. A covariate-dependent generalization of the Davidson tie model
\cite{davidson1970ties} uses
\begin{align}
D_{jico}&=e^{\eta_{jico}/2}+e^{-\eta_{jico}/2}
+e^{\kappa_{jico}},\\
P_j(Y=1)&=e^{\eta_{jico}/2}/D_{jico},\\
P_j(Y=-1)&=e^{-\eta_{jico}/2}/D_{jico},\\
P_j(Y=0)&=e^{\kappa_{jico}}/D_{jico}.
\end{align}
One explicit working predictor is
\begin{align}
\eta_{jico}&=\alpha_j+d_j(\beta+b_i+b_{iv}+b_{ic})
+\gamma_j^Oo+\gamma_j^L\ell_{ic}\nonumber\\
&\quad+\gamma_j^Ss_{jic}+u_{j,d(i)}+f_{\mathrm{fam}(j),ic},\\
\kappa_{jico}&=\kappa_j+t_i+t_{iv}+z_{ic}^{\top}\xi_j,
\end{align}
with $d_j>0$ and partially pooled effects. Identifiability constraints set
weighted human-panel offset to zero and weighted log-discrimination to zero,
center item and variant effects, and require a connected judge--cell overlap
graph with both orders and estimable covariate contrasts. Family effects model
some shared jury error; conditional independence remains only a working
assumption. The two recorded tie subtypes are collapsed to zero before this
three-category model and analyzed separately.

The human scale is anchored by the frozen weighted panel rather than assuming
jury consistency equals validity. Length and family coefficients remain
descriptive unless randomized interventions identify them: removing a length
association can remove genuine quality as well as bias.

\paragraph{Decision-targeted acquisition.}
Let terminal rule $a_T(\mathcal F)$ be the confidence-set rule in
\cref{sec:statistical-protocol}. The protocol fixes the normalized loss matrix
below, with rows $(A,B,E,\mathrm{NE})$ and true states $(A,B,E)$:
\begin{equation}
L=\begin{pmatrix}
0&4&2\\4&0&2\\2&2&0\\1&1&1
\end{pmatrix}.
\label{eq:loss}
\end{equation}
Stakeholder-specific alternatives are sensitivity analyses, not post-hoc
replacements. The working posterior defines
$R(\mathcal F)=\mathbb E[L\{a_T(\mathcal F),S(\Delta_F)\}\mid\mathcal F]$.
For target unit $x$ with annotation cost $k_x$, one-step expected value of
information is
\begin{align}
\operatorname{EVI}_{t-1}(x)
&=R(\mathcal F_{t-1})\nonumber\\
&\quad-
\sum_{y\in\{-1,0,1\}}\widehat p_{t-1}(y\mid x)
R\{\mathcal F_{t-1}\cup(x,y)\}.
\end{align}
Using predeclared temperature $\lambda_E$, define
\begin{align}
g_t(x)&=\lambda_E\operatorname{EVI}_{t-1}(x)/k_x,\\
r_t(x)&=\exp\{g_t(x)\}/\sum_{x'}\exp\{g_t(x')\},\\
p_t(x)&=\epsilon\pi(x)+(1-\epsilon)r_t(x),
\qquad \epsilon>0.
\label{eq:explore}
\end{align}
This gives $p_t(x)\ge\epsilon\pi(x)$ and
$\pi(x)/p_t(x)\le1/\epsilon$. EVI is myopically optimal only under a correct
posterior and no exploration constraint; we claim neither finite-horizon
optimality nor a general regret bound.

\section{Validation}
\label{sec:validation}

\begin{proposition}[Design unbiasedness]
\label{prop:unbiased}
Assume the finite frame and fixed budget $T$ are prespecified; $p_t$ and
$m_{t-1}$ are measurable before outcome $t$; sampling is with replacement from
the logged $p_t$; and \cref{eq:explore} holds. Then
\begin{equation}
\mathbb E[\psi_t\mid\mathcal F_{t-1}]=\Delta_F,\qquad
\mathbb E[\widehat\Delta_T]=\Delta_F,
\end{equation}
for any misspecified working jury model or outcome predictor.
\end{proposition}

The proof follows by conditioning on $X_t$ and appears in
\cref{app:datasheet}. Forward or prequential fitting is essential: each
prediction and sampling probability is frozen before its outcome. Ordinary
random $K$-fold fitting after adaptive collection can leak information through
future selections.

\begin{theorem}[Fixed-budget inference]
\label{thm:clt}
Under the assumptions of \cref{prop:unbiased}, let
$D_{Tt}=\psi_{Tt}-\Delta_F$. Suppose
\begin{equation}
\frac1T\sum_{t=1}^T
\mathbb E[D_{Tt}^2\mid\mathcal F_{t-1}]
\xrightarrow{p}V\in(0,\infty)
\end{equation}
for deterministic $V$, and for every $\eta>0$ the martingale Lindeberg
quantity
\begin{equation}
\frac1T\sum_{t=1}^T
\mathbb E[D_{Tt}^2\mathbf 1\{|D_{Tt}|>\eta\sqrt T\}
\mid\mathcal F_{t-1}]
\xrightarrow{p}0.
\end{equation}
Then
\begin{equation}
\sqrt T(\widehat\Delta_T-\Delta_F)
\overset{d}{\longrightarrow}\mathcal N(0,V).
\end{equation}
\end{theorem}

This is the standard martingale central-limit route
\cite{hall1980martingale}, specialized to predictable adaptive sampling.
Because $Y,m\in[-1,1]$ and $\pi/p\le1/\epsilon$, increments are bounded when
$\epsilon$ is fixed. Under an additional squared-increment law of large
numbers,
\begin{equation}
\widehat V_T=\frac1{T-1}\sum_{t=1}^T
(\psi_t-\widehat\Delta_T)^2,\qquad
\widehat{\mathrm{se}}=\sqrt{\widehat V_T/T}
\end{equation}
gives studentized fixed-budget inference. This is pointwise asymptotic, not
exact finite-sample, coverage. The protocol fixes $T$ before the first human
outcome and forbids threshold-based optional stopping.

\subsection{Contamination and Gaming Resistance}

Public static items are checked for exact, normalized, paraphrase, and
template-level overlap with model-accessible corpora where those corpora are
available. Ordering-based tests provide evidence for only particular forms of
memorization \cite{oren2024contamination}; absence of a signal is not proof of
cleanliness. A fresh private wave and generated variants reduce exposure but
create a new population, so longitudinal analyses retain bridge items and
separate within-wave model difference from between-wave drift.

Judge prompts and identities are hidden from evaluated systems. Position is
randomized, response identifiers are stripped, and controlled style/length
transformations probe judge sensitivity. Model outputs containing instructions
to the evaluator are separately flagged. The human-only sealed audit remains
the final check against jury manipulation.

\subsection{Statistical Protocol}
\label{sec:statistical-protocol}

The primary interval is two-sided for $\Delta_F$. The terminal action used by
both \cref{eq:loss} and the final report is
\begin{equation}
\widehat a=
\begin{cases}
A,&L_T>\delta,\\
B,&U_T<-\delta,\\
E,&[L_T,U_T]\subseteq[-\delta,\delta],\\
\mathrm{NE},&\text{otherwise}.
\end{cases}
\end{equation}
If $[L_T,U_T]$ has $1-\alpha$ coverage, the probability of a wrong decisive
or equivalence action is at most $\alpha+o(1)$ for this single comparison.
This does not bound
abstention, expected stakeholder loss, or multi-model ranking error.

\Cref{thm:clt} is conditional on the full realized frame, including the enrolled
rater panel and fixed response pairs. For a separate deployment-population
target $\Delta_{\mathrm{pop}}$,
\begin{equation}
\widehat\Delta_T-\Delta_{\mathrm{pop}}
=(\widehat\Delta_T-\Delta_F)+(\Delta_F-\Delta_{\mathrm{pop}}).
\end{equation}
The martingale variance covers only the first term. The second requires the
actual first-stage item and rater recruitment design and its replicate weights
or a justified superpopulation model; prompt variants cannot be bootstrapped
as independent root items.

\section{Empirical Study}
\label{sec:results}

Simulation precedes human collection. We vary $\Delta_F$ at
$\{-\delta,0,\delta\}$ and local alternatives
$\delta\pm c/\sqrt T$, tie prevalence, item correlation, judge-family
correlation, repeated-rater effects, order/length/family effects, outcome-model
misspecification, rare strata, small propensities, nonresponse, and temporal
shift. For each condition, 2,000 Monte Carlo trials compare target-random
human sampling, static independent calibration, PPAT-style residual
acquisition, variance-optimal residual sampling, pure uncertainty sampling,
and EVI. Each acquisition policy is crossed with the same valid
\cref{eq:estimator}; augmentation is then ablated at fixed policy. Primary
metrics are bias, RMSE, 95\% coverage, interval width, false decisive actions,
equivalence decisions, abstention, realized loss, maximum weight, and ESS,
with Monte Carlo uncertainty.

The planned real study freezes 600 probability-weighted root items spanning
declared product domains, two validated prompt forms per item, and paired
response draws from two fixed open model revisions. Four model-judge families
score all cells in both orders. A 200-rating uniform human pilot precedes 800
sequential ratings; a separate 200-rating target-random audit sampled from
\cref{eq:target-mass} remains sealed. All counts and fixed budget $T=1000$ are
set by simulation before the first human outcome. A frozen panel of 40 raters
is recruited against a written definition, assigned target weights, trained on
a fixed handbook, screened with non-study examples, and capped at 35 study
ratings in randomized task order. Identifiers are hidden, both display orders
are balanced in expectation, and every skip or nonresponse is retained.

\subsection{Main Findings}

\begin{table}[H]
\caption{Primary finite-frame decision results. Dashes indicate that human
collection has not been run. The table will report tie rate and interval, not a
leaderboard rank.}
\label{tab:main}
\centering
\begin{small}
\setlength{\tabcolsep}{2.5pt}
\begin{tabular}{lcccc}
\toprule
Design & $n_H$ & $\widehat\Delta_F$ [95\% CI]
& Tie & Action \\
\midrule
Target random & -- & -- & -- & pending \\
Static calibration & -- & -- & -- & pending \\
PPAT-style & -- & -- & -- & pending \\
EVI & -- & -- & -- & pending \\
EVI + SAHH & -- & -- & -- & pending \\
Sealed audit & -- & -- & -- & pending \\
\bottomrule
\end{tabular}
\end{small}
\end{table}

Efficiency is claimed only if interval validity is retained in simulation and
an interval for the adaptive-minus-audit difference lies within a predeclared
equivalence margin, accounting for any shared-rater covariance. Jury-only
rankings are diagnostic and never substitute for \cref{eq:estimand}.

\subsection{Robustness and Ablations}

We ablate augmentation, EVI, the exploration floor, tie modeling, root-item
effects, judge-family effects, and prequential fitting while crossing each
acquisition policy with a common estimator. Acquisition trained on
older model pairs is tested on a held-out stronger pair to expose frontier
shift. Additional falsification tests swap AB/BA order, control response length
and style, hold out prompt families, remove one judge family, change the
practical margin, and repeat the study on a fresh temporal wave.

Model checks use held-out log score, multiclass Brier score, tie calibration,
swap symmetry, posterior predictive checks, and residuals by domain, order,
length, family, judge, and item. Disagreement cases are sampled by a
prespecified stratified rule rather than selected for narrative convenience.

\section{Limitations}
\label{sec:limitations}

The finite frame may omit important deployment requests, and a convenience
rater panel supports inference only to that panel. Human preference is not
factual correctness, safety, fairness, or welfare. The Davidson model can be
misspecified, judge correlations can be more complex than modeled, and
non-transitivity can make broader rankings incoherent. Design unbiasedness does
not guarantee useful finite-sample intervals when exploration probabilities
are small.

The method requires exact logging of every adaptive probability and a fixed
budget or sequentially valid interval. EVI may save no labels and can be worse
than uniform sampling under model shift. A sealed audit consumes additional
human effort. Fresh items reduce exposure but impair longitudinal
comparability, while black-box contamination cannot be ruled out. Finally, one
pairwise replacement decision does not establish a general benchmark or model
ordering.

\section{Conclusion}
\label{sec:conclusion}

We framed language-model comparison as a decision about a frozen, weighted
human panel rather than a latent jury score. Randomized exploration and
draw-level sequential augmentation permit outcome-adaptive human acquisition
without surrendering design identification under explicit assumptions. The
formal result does not show that acquisition is efficient or that the working
model is valid. The decisive next evidence is coverage under misspecified,
correlated simulations followed by agreement with a sealed uniform human
audit.

% Acknowledgements must not appear in the anonymous review submission.
% \section*{Acknowledgements}

\section*{Impact Statement}

Evaluation methods shape which systems are deployed and which capabilities are
optimized. A narrow or unrepresentative item and rater frame can legitimize a
harmful deployment even when its confidence interval is statistically valid.
Adaptive labeling may undersample rare communities unless target weights and
the exploration floor preserve them. Raters may encounter disturbing content;
they require informed consent, task-specific warnings, the ability to skip
without penalty, fair compensation, and access to support. Deployment logs
must be de-identified and retained under strict access controls. Public
reporting will include the population definition, tie rate, abstentions,
maximum sampling weights, subgroup support, and unresolved validity threats.
The method is not intended to turn human preference into a universal quality
score or to justify deployment without separate safety and factuality audits.

\bibliography{references}
\bibliographystyle{../template/icml2026/icml2026}

\clearpage
\appendix
\onecolumn

\section{Datasheets and Reproducibility Details}
\label{app:datasheet}

\subsection{Required Evaluation Manifest}
The artifact records the frame timestamp, root-item source and target weight,
deduplication cluster, domain and language, every prompt variant and weight,
response coupling and seed, model and tokenizer revisions, decoder parameters,
judge model and prompt revisions, order randomization, rater eligibility and
training, rubric version, tie subtype, annotation cost, all $p_t(x)$ values
before outcome reveal, outcome-model revision, EVI values, stopping
rule, and exclusion reason. Probability normalization and positivity are
checked online; a failed check halts collection.

\subsection{Proof of \cref{prop:unbiased}}
Conditioning on pre-outcome history,
\begin{align}
\mathbb E[\psi_t\mid\mathcal F_{t-1}]
&=\sum_xp_t(x)\left[
\bar m_{t-1}+\frac{\pi(x)}{p_t(x)}
\{\mu_x-m_{t-1}(x)\}\right]\\
&=\bar m_{t-1}+\sum_x\pi(x)\{\mu_x-m_{t-1}(x)\}
=\Delta_F.
\end{align}
Averaging over fixed $T$ proves the second statement. Neither step requires a
correct jury model.

\subsection{Human Audit}
The audit draws complete $(i,c,h,o)$ units from $\pi$ before adaptation and is
stored behind a separate random seed. It estimates the same finite-frame target
by the unweighted sample mean; if operational assignment changes its draw law,
the exact target-to-draw ratio is retained.
Its outcomes are opened only after the adaptive algorithm, model
hyperparameters, and primary analysis are frozen.

\section{Additional Validity Evidence}
\label{app:additional}

\begin{table}[h]
\caption{Mandatory validity report; all values are pending collection.}
\centering
\begin{tabular}{lc}
\toprule
Diagnostic & Value \\
\midrule
Simulation coverage at $\Delta_F=\delta$ & -- \\
False decisive-action rate & -- \\
Maximum $\pi(X_t)/p_t(X_t)$ & -- \\
Adaptive-sample ESS & -- \\
Human inter-rater agreement & -- \\
Jury multiclass Brier score & -- \\
AB/BA swap discrepancy & -- \\
Adaptive--audit estimate difference & -- \\
Temporal-wave drift & -- \\
\bottomrule
\end{tabular}
\end{table}

Full per-domain estimates, tie subtypes, disagreement matrices, acquisition
trajectories, all simulation conditions, and randomly selected qualitative
cases will be reported. Post-hoc analyses are labeled exploratory and cannot
replace the endpoints in \cref{tab:main}.

\end{document}
